\subsection{What Migration Actually Costs}
\label{sec:online:M}

Every result above depends on $M$, the cost of moving a replica expressed in
rounds of ordinary replication, which we had assumed to be $3$. It had not been
measured, because DEFT's backup discards the state it receives: its
\texttt{update\_state} prints and returns, and the persistence code is commented
out. We therefore built the state it should have kept --- a connection-tracking
table keyed on the 5-tuple, with incremental delta replication, TCP state
tracking and idle eviction, driven by Poisson flow arrivals with truncated-Pareto
lifetimes --- and measured it (Table~\ref{tab:mrounds}).

Occupancy agrees with the Little's-law prediction within $3.3\%$ at every arrival
rate, which checks the generator rather than supplying a second result. The
measured $M$ is $23$--$41$ rounds, an order of magnitude above the assumption.

\paragraph{$M$ is a property of traffic churn, not of engineering.} It barely
moves with the wire encoding or the link speed, which is the most useful thing
here: a costlier encoding inflates the full transfer and the per-round delta by
the same factor, so the ratio is nearly unchanged. $M \approx 1/\delta$ for a
delta share $\delta$ of the table per round, here $1/3.3\% \approx 30$. No choice
of wire format or NIC removes it.

Two properties of the shipped encoding are worth recording even though they do
not change $M$. XML-RPC costs $7.31\times$ a packed binary layout ($321.8$ against
$44.0$ bytes per entry), and its \texttt{<int>} is a signed 32-bit integer, so an
IPv4 address at or above \texttt{128.0.0.0} overflows and must be sent as text.
A NAT mapping in the upper half of the address space cannot be transmitted as an
integer at all.

\paragraph{What this corrects.} At $M = 3$ migrating away from a degraded backup
looked cheap. At $M = 23$--$41$ the round trip, leaving and returning when the
host recovers, costs $2M$, and degradation must persist for a long time before it
pays. This revises the conclusion rather than the measurement, and it is part of
why degradation-aware migration adds so little over plain failover.
